\section{PASSIVE MITIGATION ANALYSIS}

This section presents the numerical results obtained using the Finite Elements framework developed for the analysis of periodic Euler-Bernoulli beams. The study focuses on how geometric periodicity and segment saturation parameters influence the dynamic behavior, particularly the distribution of natural frequencies and the formation of band gaps.

Numerical simulations were performed for both stepped and bi-material beam configurations via FEM, considering hinged-hinged boundary conditions. Table \ref{tab:sim1_values} shows the geometric and material properties adopted for the TLP simulations \cite{Domagalski2021}. Namely, the stepped beam configuration consists of square section segments with constant height and varying widths, while material properties are kept homogeneous. The bi-material beam configuration consists of square section segments with equal dimentions, and varying material properties.

% Stepped and bimat config -------------
\begin{table}[htbp]
\centering
\caption{Stepped and bimaterial configurations for Tension Leg Platform.}
\label{tab:sim1_values}
\begin{tabular}{ccccc}
\toprule
\multicolumn{5}{c}{Bi-material Beam} \\
\cmidrule(r){1-5}
Segment & Height ($h$) & Width ($b$) & Young Modulus ($E$) & Density ($\rho$) \\
\midrule
Middle Segment ($k$)         & $0.01$ m & $0.0250$ m & $205$ GPa & $7850$ kg/m$^3$ \\
Outter Segments ($k\pm1$) & $0.01$ m & $0.0250$ m & $69$ GPa  & $2700$ kg/m$^3$ \\
\midrule
\multicolumn{5}{c}{Stepped Beam} \\
\cmidrule(r){1-5}
Segment & Height ($h$) & Width ($b$) & Young Modulus ($E$) & Density ($\rho$) \\
\midrule
Middle Segment ($k$)         & $0.01$ m & $0.0375$ m & $205$ GPa & $7850$ kg/m$^3$ \\
Outter Segments ($k\pm1$) & $0.01$ m & $0.0250$ m & $205$ GPa & $7850$ kg/m$^3$ \\
\bottomrule
\end{tabular}
\end{table}
% --------------------------------------

% The first 10 natural frequencies are calculated for a fixed saturation parameter $\alpha=1/2$, which defines the relative length of the central segment within each unit cell. A comparison between analytical and FEM results is presented in Table \ref{tab:sim1}. The two models show good overall agreement and maximum percentil error of $0.0006029 \%$ for the bi-material configuration, and $0.000688 \%$ for the stepped configuration. There is a notable increase in percentil error with the increased mode number. A similar analysis made by \textcite{Domagalski2021} yielded resembling results.

The first 10 natural frequencies are calculated for a fixed saturation parameter $\alpha=1/2$, which defines the relative length of the central segment within each unit cell. A comparison between analytical and FEM results is presented in Table \ref{tab:sim1}. The two models show good agreement and maximum percentil error of 0.0006029\% for the bi-material configuration, and 0.000688\% for the stepped configuration. There is a notable increase in percentil error with increased mode numbers. A similar analysis by \textcite{Domagalski2021} yielded similar results.
% Frequências Naturais + erro --------
\begin{table}[htbp]
\centering
\caption{Comparison of the first 10 natural frequencies (in rad/s) for stepped and bi-material configurations, with $\alpha=1/2$.}
\label{tab:freq_validation}
\begin{tabular}{c ccc ccc}
\toprule
\multirow{2}{*}{Mode} & \multicolumn{3}{c}{Bi-material Beam} & \multicolumn{3}{c}{Stepped Beam} \\
\cmidrule(lr){2-4} \cmidrule(lr){5-7}
 & Analytic & FEM & Error (\%) & Analytic & FEM & Erro (\%) \\
\midrule
1 & 315.2802 & 315.2802 & 0.0000026 & 404.4794 & 404.4794 & 0.0000014 \\
2 & 1263.0599 & 1263.0600 & 0.0000028 & 1619.0215 & 1619.0216 & 0.0000036 \\
3 & 2849.4788 & 2849.4789 & 0.0000031 & 3647.0420 & 3647.0422 & 0.0000048 \\
4 & 5085.7799 & 5085.7807 & 0.0000149 & 6494.4482 & 6494.4494 & 0.0000186 \\
5 & 7989.5603 & 7989.5633 & 0.0000375 & 10169.3888 & 10169.3935 & 0.0000465 \\
6 & 11585.8390 & 11585.8480 & 0.0000779 & 14680.1879 & 14680.2025 & 0.0000994 \\
7 & 15906.9949 & 15907.0185 & 0.0001480 & 20027.2895 & 20027.3276 & 0.0001903 \\
8 & 20990.2931 & 20990.3477 & 0.0002602 & 26171.8004 & 26171.8893 & 0.0003397 \\
9 & 26871.9545 & 26872.0700 & 0.0004297 & 32855.4483 & 32855.6399 & 0.0005831 \\
10 & 33679.3825 & 33679.5855 & 0.0006029 & 47348.2847 & 47348.6104 & 0.0006880 \\
\bottomrule
\end{tabular}
% \\ \vspace{2pt} \footnotesize{* ... .}
\label{tab:sim1} 
\end{table}
% ---------------------------------------

To characterize the periodic effects on the TLP structure, the first $22$ natural frequencies were computed for a fixed saturation parameter $\alpha=1/2$. For the bi-material configuration, a single Bragg scattering band gap ($\Delta\omega=34422.19$ rad/s) was detected between modes $20$ and $21$ (Figure \ref{fig:sim1}a). Analysis of the stepped beam configuration revealed two distinct gaps: one between modes $9$ and $10$ ($\Delta\omega=14493.25$ rad/s), and a second between modes $20$ and $21$ ($\Delta\omega=34728.40$ rad/s) (Figure \ref{fig:sim1}b).
%--------------------------------------
\begin{figure}[htbp]
    \centering
    \caption{Band gaps for (a) bimaterial and (b) stepped beam configuration, with $\alpha=1/2$.}
    \includegraphics[scale=.53]{bandgap-bimat-8_16.png}
    \includegraphics[scale=.53]{bandgap-stepped-8_16.png}
    {\\ \hfill (a) \hfill (b) \hfill } %
    \label{fig:sim1}
\end{figure}
%----------------------------------------

Changes in the saturation parameter $\alpha$ lead to noticeable shifts and broadening of bang-gap positions over the frequency spectrum. This result can be seen on Figures \ref{fig:sim2} and \ref{fig:sim3}. The former shows evolution of band gaps in natural frequencies between modes number $7$ and $11$, for $\alpha=i/16$, where $i=1,2,\cdots,15$. The latter shows the same for modes $18$ through $22$.
% Saturation parameter (modes 7 through 11)
%--------------------------------------
\begin{figure}[htbp]
    \centering
    \caption{Frequency gaps broadening by saturation parameter $\alpha$ for (a) bi-material and (b) stepped beams, for natural modes $7$ through $11$.}
    \includegraphics[scale=.53]{natfreq_7-11_bimaterial.png}
    \includegraphics[scale=.53]{natfreq_7-11_stepped.png}
    {\\ \hfill (a) \hfill (b) \hfill } %
    \label{fig:sim2}
\end{figure}
%----------------------------------------

% Formation of band gaps can be evaluated on Figure \ref{fig:sim2}a for bi-material configuration, by observation of vertical space between natural frequency values for various modes, for the same saturation parameter ($\alpha$). For modes number $7$ through $11$, natural frequency varies considerably between modes $10$ and $11$, with a point of maximum on $\alpha=3/16$ ($\Delta\omega=9363.2$ rad/s), which was also the highest variation on lower frequencies. Beyond $\alpha=0.5$, the variation takes place between modes number $9$ and $10$, and exhibits similar behavior, with a point of maximum on $\alpha=13/16$ ($\Delta\omega=8846.2$ rad/s). However, this behavior does not result in formation of band gaps in lower frequencies, as can be seen by comparison to prismatic frequencies on Figures \ref{fig:sim1} and \ref{fig:sim2.check.bimat} for saturation parameters $\alpha=0.5$, $\alpha=3/16$ and $\alpha=13/16$.

Band gap formation in the bi-material configuration is analyzed in Figure \ref{fig:sim2}a via the spectral spacing between modes for varying $\alpha$. In the range of modes $7$ to $11$, the frequency gap peaks between modes $10$ and $11$ at $\alpha=3/16$ ($\Delta\omega=9363.2$ rad/s). For $\alpha>8/16$, the primary gap shifts to modes $9$ and $10$, peaking at $\alpha=13/16$ ($\Delta\omega=8846.2$ rad/s). Despite these spectral distortions, a comparison with prismatic beam frequencies (Figures \ref{fig:sim1} and \ref{fig:sim2.check.bimat}) indicates that no effective band gaps are formed in the lower frequency spectrum.
% Bi-mat config for lower freq showing no bandgaps
%--------------------------------------
\begin{figure}[htbp]
    \centering
    \caption{Band gap absence for bi-material beam configuration on points (a) $\alpha=3/16$ and (b) $\alpha=13/16$.}
    \includegraphics[scale=.53]{bandgap-bimat-3_16.png}
    \includegraphics[scale=.53]{bandgap-bimat-13_16.png}
    {\\ \hfill (a) \hfill (b) \hfill } %
    \label{fig:sim2.check.bimat}
\end{figure}
%----------------------------------------

The stepped beam configuration exhibits a monotonically increasing frequency gap between modes $9$ and $10$ (Figure \ref{fig:sim2}b), peaking at $\Delta\omega=17044.0$ rad/s for $\alpha=11/16$. Comparative analysis with the prismatic spectrum confirms the onset of band gaps for $\alpha\geq 6/16$ (Figure \ref{fig:sim2.check.stepped}a).
% Stepped config for lower and higher freq showing bandgaps
%--------------------------------------
\begin{figure}[htbp]
    \centering
    \caption{Band gaps for stepped beam configuration for (a) $\alpha=6/16$ and (b) $\alpha=14/16$.}
    \includegraphics[scale=.53]{bandgap-stepped-6_16.png}
    \includegraphics[scale=.53]{bandgap-stepped-14_16.png}
    {\\ \hfill (a) \hfill (b) \hfill } %
    \label{fig:sim2.check.stepped}
\end{figure}
%----------------------------------------

% Analysis of mode numbers $18$ through $22$ show that, for the bi-material configuration (Figure \ref{fig:sim3}a), the structure shows minor frequency gaps between modes number $19$ and $20$, for $\alpha \in [1/10,\: 3/16]$. Which are ``absorbed'' and later amplified by modes number $20$ and $21$. For $\alpha \in [3/16,\: 8/16]$ the scattering widens in both directions of the spectra, until a maximum $\Delta\omega = 34434.56$ rad/s is reached for $\alpha=0.5$. The analysis also shows that the band gap behavior is symmetric about $\alpha=0.5$, i.e., it shortens for $\alpha\in [8/16,\: 15/16]$ before being partially absorbed by lower frequencies. As seen on Figure \ref{fig:sim2.check.bimat}, the identified frequency gaps do not translate as band gaps until saturation parameters closer to $\alpha=0.5$. Which means that the introduction of material periodicity has a tendency to only generate band gaps for limited configurations, allowing for poorer tuning capabilities.

Analysis of modes $18$ through $22$ in the bi-material configuration (Figure \ref{fig:sim3}a) reveals minor frequency gaps between modes $19$ and $20$ for $\alpha\in[1/10,3/16]$. These are subsequently shifted to modes $20$ and $21$. For $\alpha\in[3/16,8/16]$, the gap broadens spectrally, reaching a maximum $\Delta\omega=34434.56$ rad/s at $\alpha=8/16$. The behavior is symmetric about $\alpha=8/16$, narrowing for $\alpha\in[8/16,15/16]$, as the mode shapes evolve and the stopband boundaries shift relative to the discrete natural frequencies. As proved from Figure \ref{fig:sim2.check.bimat}, these spectral gaps only manifest as true band gaps for saturation parameters near $\alpha=8/16$. Thus, material periodicity shows limited capabilities for band gap generation, resulting in restricted tunability.
% Saturation parameter (modes 18 through 22)
%----------------------------------------
\begin{figure}[htbp]
    \centering
    \caption{Band gap broadening by saturation parameter for (a) bi-material and (b) stepped beams, for natural modes $18$ through $22$.}
    \includegraphics[scale=.53]{natfreq_18-22_bimaterial.png}
    \includegraphics[scale=.53]{natfreq_18-22_stepped.png}
    {\\ \hfill (a) \hfill (b) \hfill } %
    \label{fig:sim3}
\end{figure}
%----------------------------------------

% The stepped configuration presents a slightly different behavior for higher modes (Figure \ref{fig:sim3}b). The frequency gap widens from $\alpha=2/16$ onwards, reaching the absolute maximum gap ($\Delta\omega=37024$ rad/s) for $\alpha=7/16$, and between modes number $20$ and $21$. This gap also shifts location upwards along the frequency spectra for $\alpha\in [1/16,\: 10/16]$, between modes number $20$ and $21$. However, for $\alpha > 0.5$ the band gap is absorbed by lower modes, reaching a local point of maximum of $\Delta\omega=35502$ rad/s, for $\alpha=14/16$ between modes number $19$ and $20$. Careful consideration of Figures \ref{fig:sim2}, \ref{fig:sim2.check.stepped} and \ref{fig:sim3} leads to the conclusion that the introduction of geometric periodicity allows for two band gaps for varied configurations of the saturation parameter. This turns the stepped beam into a superior, more tunable structure than the bi-material beam, towards the goal of vibration control in Tension Leg Platforms.

The stepped configuration exhibits distinct behavior for higher modes (Figure \ref{fig:sim3}b). The frequency gap widens starting at $\alpha=2/16$, reaching an absolute maximum ($\Delta\omega=37024.0$ rad/s) at $\alpha=7/16$ between modes $20$ and $21$. This gap shifts spectrally upwards for $\alpha\in[1/16,10/16]$. However, for $\alpha>8/16$, the attenuation zone shifts to lower mode orders, achieving a local maximum of $\Delta\omega=35502.0$ rad/s (at $\alpha=14/16$) between modes $19$ and $20$. Analysis of Figures \ref{fig:sim2}, \ref{fig:sim2.check.stepped}, and \ref{fig:sim3} confirms that geometric periodicity induces two distinct band gaps across various saturation parameters. This establishes the stepped beam as a superior, more tunable structure compared to the bi-material configuration for vibration control in Tension Leg Platforms.